% ============ 第三部分：第 9--12 章、附录与参考文献 ============

\section{数学界的验证与约束}

前三章处理的是“构造是怎么做的”。这一章处理“别人怎么看”，并特别关注一条容易被忽略的区别：验证文献里既有\emph{否证性}的结论，也有\emph{约束性}的结论，二者对构造的意义完全不同。

\subsection{Tao 的独立评估与方法溯源}

Terence Tao 在 9 月 7 日的博文\cite{tao-blog-sep7}中给出了一份独立于两家机构的评估。他称 Alpöge--Buckmaster 的结果是“了不起的成就”，把基本路线的归属明确记在 Córdoba 与 Martínez-Zoroa 名下，并指出：这条路从“外力下爆破”延伸到 \NS{} 是原则上可预期的，甚至有机会去掉外力。值得注意的是，这篇博文写于 OpenAI 公布 \NS{} 结果之前——也就意味着，“从 Alpöge--Buckmaster 的强迫 Euler 到千禧年问题，道路对任何理解这套方法的人都是可见的”这一判断，是一份独立的、时序在先的材料。

Tao 同时强调，他仍需要更多时间来消化细节，并把“突出这些工作的核心想法”视为这类评述的主要价值——求解问题只是发展数学理解的代理目标，而非目的本身。这句话为第 10 章的认识论讨论定了调。

\subsection{Constantin--Ignatova--Vicol 的刚性定理}

9 月 17 日，Peter Constantin、Mihaela Ignatova 与 Vlad Vicol 上传了一篇预印本\cite{civ}，其结论初看像是冲着 OpenAI 的构造去的，实则需要仔细分辨方向。

他们的定理说：设 $(u,\pi)$ 是带外力的三维 \NS{} 在单位抛物柱上的适当弱解，在紧子集上光滑；设外力在奇点时刻前一致地有 $C^{2}$ 界，且在空间变量上局部一致实解析；再设速度满足两条与手稿构造相对应的性质——(i) 其角向平均 $\mathcal{P}u$ 满足各向异性的 Type II 界（$v_{r}$ 被 $\tau^{-1/2}$ 控制，$v_{z},v_{\theta}$ 被 $\tau^{-1/2-h}$ 控制，每求一次 $r$ 导数付 $\tau^{-1/2}$、每求一次 $z$ 导数付 $\tau^{-1/2+h}$），(ii) 在一个坍缩的核心球上速度严格轴对称——那么该假想奇点实际上是正则点。

换言之：在这两条性质加上解析外力成立时，奇点不可能出现。由于 OpenAI 的构造恰恰满足 (i) 与 (ii)，定理的推论就变成一句约束：任何具有性质 (i)、(ii) 且外力在奇点时刻前保持 $C^{2}$ 有界的构造，其外力既不能在奇点附近恒为零，也不能在空间变量上实解析。作者还特意在手稿的附录里核对了 OpenAI 构造确实满足 (i) 与 (ii)，并把对方的外力性质定为“$C^{\infty}$ 光滑但非实解析”。

\subsection{关键：这是约束，不是否证}

初学者容易把这条定理误读为“OpenAI 的构造被证伪了”。恰恰相反。

手稿的外力由紧支撑的截断拼出，因而是光滑但\textbf{非}解析的（第 7 章已说明：紧支撑的非零实解析函数不存在，奇点邻域也非解析）。刚性定理说的是“若外力解析则正则”，而构造的外力不解析，因此它落在定理假设之外——定理对它不作任何断言，更谈不上否定。定理真正的价值在三处：其一，它说明了构造中哪些特征是\textbf{本质}的——各向异性 Type II 标度、坍缩核心的精确轴对称、以及外力不解析，这三者缺一不可，否则奇点会被刚性逼回正则；其二，它给出了一个可判定的“必要形态清单”，为后来者提供了检查表；其三，作者明确声明“我们并不声称验证了手稿构造的正确性”，只是从手稿陈述中抽取所需性质。这种克制本身就值得记录。

\begin{principlebox}
\textbf{约束性结论的正确读法}\quad “在更强假设下结论为正则”，并不等于“原构造错误”，而是给出“原构造必须避开这些假设”这一必要条件。CIV 的贡献在于把“外力必须非解析”从技术细节升格为结构必然。
\end{principlebox}

\subsection{刚性定理的证明思路：放大尺度与古代解}

CIV 的结论不只是一条约束，它的论证方式本身值得一记，因为它示范了当代正则性理论的一种标准操作。思路分三步。

第一步是放大。利用各向异性的 Type II 界所提供的长度尺度 $\ell_{r}(t)=(-t)^{1/2}$、$\ell_{z}(t)=(-t)^{1/2-h}$，在假想奇点处做抛物型伸缩。这一步的关键在于：Type II 标度允许把速度场规范化，使伸缩后的解在紧集上一致有界——否则连取极限都无从谈起。

第二步是取“古代解”。伸缩参数趋于零时，得到的极限解定义在 $(-\infty,0]$ 上，即一个自古以来就存在的解。这类对象的约束极强：它没有初始时刻，因而任何依赖“信息从初值单向流入”的论证都被排除。

第三步是刚性。对古代解施加条件（外力的解析性、解的轴对称性），证明这样的古代解必然平凡，从而把假想奇点排除。作者指出，这一步的机制与他们此前关于轴对称 \NS{} 的工作一脉相承。

这三步合起来说明：刚性定理真正用到的是假想奇点附近解的形态足够特殊这一事实；它并未用到被质疑的引理 8.6。两条验证线互不干涉，这也是本文把它们分开处理的理由。

\subsection{Clay 的程序}

克莱数学研究所 9 月 11 日的表态\cite{clay-news}只说问题“看上去已被解决”，并强调评估与授信的流程“刻意不匆忙”。按其规则，授奖需要满足若干条件，包括在合格期刊发表、发表后至少两年、以及数学界的普遍接受；公开材料显示至今没有授奖决定，且 OpenAI 表示不打算为此申领那笔奖金。把这三件事放在一起看，可以得到一个平稳的结论：Clay 的表态是“注意到了一项重要声明”，不是“认证了一个定理”。

\subsection{形式化验证能保证什么、不能保证什么}

Lean 证书给出的是：某个命题在一套公理系统内可证，且证明项合法。它不能自动给出的是：这个形式化命题与手稿中那句自然语言命题是否\textbf{同义}。二者之间的空隙，正是第 8 章所述争议的所在。对此，学界有一句被反复引用的比方——耶胡达·拉夫（Yehuda Rav）说“定理是标题，证明是正文”：一个正确但未被消化的证明，只给了标题而没给正文。这句话出自 Tapio Schneider 发表于 Tao 博客的客座文章\cite{schneider}，它把讨论从“对错之争”引向“理解之争”，也正是下一章的主题。

\section{认识论与学术生态：形式化、理解与归责}

\subsection{形式化不等于理解}

把前几章的线索接起来，可以得到一个不依赖任何争议事实的判断：形式化验证与数学理解是两件不同的事。形式化解决的是“从公理到命题的推理是否合法”，理解解决的是“人为什么能看出这个命题应当成立、这套方法能迁移到哪里”。前者可以自动化，后者目前不能。Schneider 的文章提议把 AI 放进“可审计的脚手架”里使用——在数学里是 Lean，在物理与工程里是高分辨率模拟、观测或实验——让 AI 在每一环都可被单独检验的位置上发挥作用。他的例子很干净：天气预报这类每天都能核对的预测，黑箱模型可以独立成立；而气候外推、飞行器定型这类预测必须在被经验检验之前就获得信任，此时可审计的因果链才是信任的来源。

值得补一句的是：形式化与理解的分野，在数学内部早已有先例，并非 AI 带来的新问题。一个极端的例子是四色定理：1976 年由计算机辅助完成，其论证此后被部分学者质疑过“是否算证明”，直到后来的形式化工作把可核验的部分做得更彻底。另一个例子是有限单群分类：其证明长达上万页，几乎没有人从头到尾读完，社区接受它靠的是分层信任——大量中间结果被反复重证与引用。数学共同体早就活在“不完全理解却接受结论”的状态里；AI 只是把这种状态从“少数巨型证明”扩展到了更多场合，并让它变得廉价而频繁。

\subsection{理解的代价：为什么“消化”如此重要}

还有一层更温和但同样重要的理由，说明为什么“可理解”不只是审美要求。数学知识的再生，靠的是方法在人群之间的传播：一个被理解的构造会衍生出新的想法、新的练习、新的应用；一个只被机器检验过的构造则停在那里，除非有人把它想通。Rav 的“标题与正文”是为了说明这一点，Schneider 的 episteme/techne 之辨也是为了说明这一点。

具体到本项工作，“消化”的难度是肉眼可见的：166 页正文加两百余万行 Lean；外部读者若要完整核对，需要同时具备调和分析、流体方程与形式化工具三方面的能力。这不是 OpenAI 独有的问题——Buckmaster 也坦承自己那篇 Euler 写稿“像 AI 生成垃圾”，并为此道歉。当构造的复杂度超过共同体的消化带宽时，问题不是“对不对”，而是“多久以后才能真正被人用起来”。

\begin{quotebox}
“定理是标题，证明是正文。”——一个正确却不被理解的证明，给出了标题，却没有交代清楚证据链。Schneider 的论点，是把 AI 安排在证据链上那些能够被单独检验的环节，而不是让它端到端地给出结论。
\end{quotebox}

\subsection{episteme 与 techne}

Tao 与 Schneider 都借用了古希腊的一对概念：episteme（解释性理解）与 techne（预测与制作的技艺）。求解一个方程的解的性态，属于 episteme；让模型预报明天的天气，属于 techne。Schneider 的观察是，当预测能被廉价地反复核验时，techne 可以独立成立（黑箱天气预报就是例子）；当预测必须在被核验之前就被信任时，episteme 与 techne 便不可分离，可审计的证据链随之成为硬约束。这一段论证对数学的启示是：Lean 让它变成“可以被机器核验的 techne”，但数学共同体真正消费的仍是 episteme——这解释了为什么一个 Lean 编过的证明，仍可能被数学界抱怨“没有交付理解”。

\subsection{优先权、训练数据与署名的争议记录}

围绕本项工作产生的一组事实主张，涉及研究生态与科研伦理，而不涉及数学内容本身。以下叙述仅依据公开材料并逐一标明来源。这些主张之间存在直接冲突，本文不掌握足以裁断的材料，因而不对任何一方的主张作出裁断。

关于优先权，据 Buckmaster 的公开声明\cite{buckmaster-statement}，他与 Alpöge 在 8 月中旬已获得光滑外力下 IPM、Boussinesq 与三维 Euler 方程的爆破结果，并认为自己同时握有次耗散 \NS{} 的结果，只是由于 Lean 验证尚未完成而未一并发布。OpenAI 的官方说明\cite{openai-blog}则称，其相关工作自 9 月 1 日因听闻“两个千禧年问题已被解决”的传闻而启动；公司表示其研究人员与智能体在对方公开发布之前没有看到对方的任何工作，并说明在 9 月 6 日完成 Lean 验证后曾联系对方，提出同步发布、承认对方在外力 Euler 上的优先权。双方对若干关键时间节点的叙述存在出入，特别是首个提示词的发出时间；在缺乏版本历史等客观记录的情况下，这一争议无法由外部材料裁决。

关于训练数据，据 Buckmaster 的声明\cite{buckmaster-statement}，他曾询问其与 Alpöge 存放在 Codex 会话中的草稿是否被用于训练或检索，未获得直接答复。OpenAI 方面表示，其模型没有查询特定用户数据，同时承认“虽不太可能，但不能排除来自用户使用的去标识化数据对模型改进有所贡献”。围绕这一表述，后续的技术讨论提出了一种更窄、更可检验的解读：相关机制更可能是对未选择退出的用户对话进行强化学习，而非直接“查阅草稿”；真正待定的问题因而变为最近一次训练发生的时间，以及模型是否会记住强度足以影响结论的片段。这一表述与 OpenAI 的措辞相容——公司否认访问特定用户数据，但不排除去标识化的派生数据对模型有所贡献。

关于署名，手稿封面唯一的署名是“OPENAI”，这一机构署名不区分研究纲领的来源、中间推进、\NS{} 特有的构造与形式化工作。对此的批评集中在责任与归责上：当外部研究者质疑某个引理或形式化对应关系时，由谁出面回应并不明确。支持机构署名的常见理由则是：成果确系组织内部的系统与团队共同产生，个体贡献难以按传统学术出版的方式切分。这一分歧是学术出版规范在面对新型研究组织形态时的一个具体案例。

关于通话内容的争议，据 Buckmaster 公布的叙述，他在 9 月 6 日与 OpenAI 方的通话中被提出两种发表方案，并记录了若干他本人视为施压的话语。被点名的 OpenAI 方面研究者公开称这些指控“虚假且具有煽动性”，表示自己是遵循学术规范参与讨论，并承诺给出更完整的说明。同一事件的两种叙述就此直接冲突。

需要区分的是“数学有效性”与“产生条件是否正当”这两件事：一个证明可以在产生条件存疑的情况下依然正确；反过来，产生条件的争议也不构成对该定理本身的数学反驳。第 9 章的约束性结果属于数学层面；本节的记录属于事实主张层面，且这些主张的来源、口径与可核验程度各不相同，读者宜按上文的归因分别看待。

\subsection{七百余篇手稿与负责任披露之争}

10 月的报道称，OpenAI 已放出数百篇 AI 生成的数学手稿，并与学界就披露方式产生摩擦：一些参与过 8 月咨询的数学家反映，他们希望以正式论文而非博客或社交媒体形式披露的意见未被采纳。把这些放进本文的框架里看，它是第 10 章那组张力的一次放大——当产出速度远超共同体消化速度时，验证、归属与理解三件事会同时承压。本节仅记录事实，不作价值判断。

\subsection{小结}

把这一章压成一句：\NS{} 事件的数学部分是清晰的（一条被验证过方向、被约束过形态的构造路线），它的社会学部分则远未尘埃落定。对本文的读者而言，重要的能力是能同时握住这两条线，而不把它们混为一谈。

\section{诚实的边界与开放问题}

\subsection{备选 A、B 仍未决}

所有公开材料一致指出：本次结果是备选 C 与 D，无外力的 A 与 B 依然开放。9 月 20 日，Petrillo 与 Glimm\cite{petrillo-glimm}提出了“正缺陷问题”作为无外力搜索的靶标：要求周期盒上从光滑数据出发的 Leray--Hopf 解在某个有限时间窗内损失的能量，超出黏性所耗散的部分。正缺陷蕴含爆破，因而蕴含备选 B 的否定；其逆向是否成立尚不清楚。他们把该靶标等价地化为 Littlewood--Paley 壳层能量通量的下界（即 Onsager 湍流理论中粗粒化通量的 Fourier 侧形式在窗口上的平均），并列出候选解必须满足的必要条件。这项工作本身也在 Lean 4 中形式化，且刻意在假设里引入 \NS{} 对象，使方程只通过假设进入。它对本文的意义是：它把“无外力路径到哪一步、缺什么”讲清楚了——缺的不是方法，而是一个靶标。

\subsection{解析外力障碍：可绕开还是本质}

第 9 章的刚性定理把外力非解析变成了必要条件。接下来的问题是：这一障碍是构造的偶然特征，还是 \NS{} 强迫爆破的固有属性。若它是固有属性，那么“光滑外力”这条路本身就带着一条暗线约束——外力必须在奇点附近保持非解析；若它只是当前构造的副产品，未来或许能构造出外力更“规整”的爆破。公开材料目前不足以回答，但它给出了一个明确的检验方向：尝试在保持 (i)(ii) 的前提下改进外力的解析性，看刚性定理在哪一点被触发。

\subsection{引理 8.6 落差的影响范围}

如第 8.5 节所述，这一落差是否传导到主定理，取决于该引理在修正链中的调用结构与余量。落差的“Lean 一侧”已无争议：本文据公开仓库 commit \texttt{f9e8bc5} 一手读取了该声明（文件 \texttt{NavierStokes/SmoothFamilyTorusInverse.lean} 第 1059 行，两个假设项取 \texttt{w.length + 5}，权重指数为 4，详见第 8.3 节），确认其字面强度确为 $m+5$；于是争点不在“Lean 写了什么”，而在“$m+4$ 的正文声明是否成立、以及在何处被调用”。一个可操作的收敛问题是：把引理 8.6 从 $m+4$ 改为 $m+5$ 后，其后各章中依赖它的估计是否仍成立；若成立，则落差是局部可修的；若不成立，则需在更上游重做。这需要逐条比对正文与 Lean 的对应关系，属于可以完成但颇为费时的核查，也是验证工作最应当优先投入的地方。

\subsection{构造能否被简化到人类可消化}

166 页正文加两百余万行 Lean，其规模本身就是理解的门槛。历史上，重大构造通常要经过数年乃至数十年的简化、重讲，才会进入教科书。是否已有或将会出现一份把核心思想压缩到数十页、可被研究生读懂的阐述，是判断这项结果能否真正进入数学共同体的关键。已有的“友好读本”（如 Shavit 的机制笔记\cite{shavit}）是这类努力的雏形，但它们解释机制，尚未重构证明。

\subsection{反方证据：把“解决”读成“终局”会错在哪里}

一篇诚实的技术综述必须把反方证据摊开。反对把此事称为“千禧年问题已解决”的理由至少有六条，它们彼此独立，且都建立在公开材料之上。

第一，方向性。结果落入备选 C、D（有外力），而最常被公众理解为“水的方程”的备选 A、B（无外力）仍然开放。用“有外力则可爆破”去回答“无外力是否永远光滑”，是换了一个问题。

第二，验证层次。Lean 通过只保证“形式命题可证”，不保证“形式命题与原命题同义”。Bastounis 等人已在引理 8.6 处给出一个具体的不一致样本。真正的独立复核（由第三方重做关键估计）尚未完成。

第三，传出规模。手稿 166 页、Lean 两百余万行，共同体的消化需要时间；Clay 的规则也明确要求“两年”与“普遍接受”。9 月的表态不是认证。

第四，构造的人为性。这是一条被精心设计的路径，依赖大量参数选择；它不给出“一般流会怎样”的刻画，因而对理解真实湍流（Kolmogorov 尺度以上的、普遍存在的湍流）帮助有限——Schneider 的文章专门指出了这一点：该结果“对理解物理湍流几乎没有发言权”。

第五，归属与条件争议。优先级、训练数据、署名方式目前均未定论，且两方叙述直接冲突。这不是数学错误，但足以说明“共同体的接受”还没有发生。

第六，机制而非结论的进展。CIV 的刚性定理表明，任何满足同样形态的构造都必须使用非解析外力。这既是重要洞察，也提示：这条路径可能带有一条内在的形态约束，而不是一条通用的爆破机制。

\begin{keybox}
\textbf{平衡读法}\quad 把这些反方理由与第 13 章的四点定位并置，得到的不是“模棱两可”，而是一个精确坐标：数学结论在特定范围内成立且已给出形式化证书（自评）；无外力情形未决；独立验证与共同体接受尚在途；物理意义有限；归属未定。把任何一项单独放大，都会得到失真的结论。
\end{keybox}

\subsection{可检验的收敛性建议：尚未完成的严格推演}

在给出方向建议之前，需要把“尚未完成”这件事说准。本项工作并非没有形式化——OpenAI 公开了包含六十余万行 Lean、三万余条定理与引理的仓库，且未留未证明的占位、未引入标准三公理之外的公理。真正尚未完成的是三件事：其一，该证书的\textbf{独立核验}（仓库自述评审状态为“自评”，未见第三方逐条复核的公开结果）；其二，证书与正文之间的\textbf{语义忠实性}（引理 8.6 处已有具体反例）；其三，其\textbf{数学思想的严格重述与推演}——即由共同体独立地、以可读的方式把“为什么这套构造能成立”重新推导一遍。第三件事之所以关键，是因为在上述两项完成之前，我们尚不能确定自己讨论的究竟是一个已确立的定理，还是一个有待确立的候选定理；而第 12 章关于“思想定位”的判断，恰恰以“我们确实理解了证明的动作”为前提。

以下任务按“先证成、再消化、后拓边”的顺序排布，供后来者按序推进。

其一，\textbf{核对引理 8.6}。对照手稿正文第 96 页附近的 (8.19) 与 Lean 仓库 commit \texttt{f9e8bc5} 中 \texttt{norm\_derivativeWord\_inverse\_le} 的假设，判定字面强度（$m+4$）是否为真；若为真，补一份与正文严格对应的形式化陈述；若为假，向上游追踪该引理被调用的位置，判断主定理是否仍成立。

其二，\textbf{复核刚性定理与技术细节的对应}。逐条核对 CIV 附录 A 中从手稿陈述抽取的性质 (i)(ii)，确认它们确实是手稿所证，而非对手稿的额外假设。

其三，\textbf{语义忠实形式化的协议}。对每一个承重引理，独立撰写一份与正文一一对应的形式化陈述，并附“对应性证据”（逐行或逐式对照说明），而不是依赖一次性把整篇论文译成代码。这是把第 8、9 章所述断层缩小的最直接手段。

其四，\textbf{构造的可简化}。尝试把“脉冲供给 Reynolds 应力”这一机制从 166 页的估计中剥离出来，写一份数十页、可被研究生独立读通的重述。

其五，\textbf{加强结论}。审视备选 C 是否可加强到“外力实解析”或“外力可取零”的变体；这一步直接与 CIV 的障碍对话，也是通往备选 A、B 的必经关口。

其六，\textbf{机制的可迁移性}。把“零平均振荡经二次自作用生成指定平均通量”抽成一个可复用的模块，测试其能否迁移到磁流体方程、次耗散 \NS{} 或其他流体模型；能否迁移，是判断它属“纲领内一步棋”还是“可推广的新工具”的实测标准。

其七，\textbf{结构分类}。以 CIV 型刚性定理为工具，刻画“哪些（标度、对称性、解析性）组合允许爆破、哪些被刚性排除”，把当前构造放进一张完整的分类表中；这既是验证工作，也可能产出独立的新定理。

其八，\textbf{无外力靶标的检验}。检验“正缺陷问题”作为无外力靶标的可行性：Petrillo--Glimm 已把靶标化为 Littlewood--Paley 壳层通量的下界，并做了 $128^{3}$、$256^{3}$ 的伪谱搜索；需要进一步判断该路径能否导向证明。

\subsection{这件事与“AI 能不能做数学”是两问}

最容易发生的误读，是把本项工作直接读成对“AI 是否会做数学”的回答。这两件事需要分开。

严格地说，本项工作证明的是一个数学命题（在其成立的前提下），而不是一个关于 AI 能力的定理。三点理由。其一，过程有大量人类参与：OpenAI 的团队设计了多智能体框架、划分了问题变体、在过程中更新模型、用 Codex 汇总各组洞见；按其自述，问题变体（A/B/C/D 分发给不同组）这一步是人安排的。其二，路线有明确的人为先驱：Córdoba--Martínez-Zoroa 的纲领、Alpöge--Buckmaster 的推进，都是人类在 AI 之前数月乃至数年铺设的。其三，验证环节依赖的是 Lean 内核，一种与人无关但对“理解”也无关的形式工具。

反过来说，也不能因为“有人类参与”就否认这次跃迁。把一份“历史上最糟糕的写稿”在七天内变成 Lean 可验证的证明，这件事本身在几年前是不可想象的；它确实标记了工具能力的台阶。诚实的立场是：这是一次“人设定方向与判据、机器承担搜索与记账”的协作突破，其意义在于改变了求解环节的成本结构，而不在于宣告某个关于智能的命题。

\subsection{长期含义}

把视野拉长，这件事提出的问题比它回答的问题更多。当“求解”的速度远超“消化”的速度时，数学共同体的训练方式、评审方式、归功方式都面临重估。Tao 的立场是：理解与求解是光谱的两端，有人以求知为目的，有人以求解为目的，二者都值得保留，不应互相规定何者为“正”。Gowers\cite{gowers}则更直接地把这称为一次“苦涩的教训”——他所在的自动定理证明小组试图通过理解人类如何证题来设计算法，而大模型绕开了这种理解照样做到了。这一判断与本文的技术观察一致：工具层面的自动化已经发生，而解释层面的自动化仍未到来。

\section{数学思想的定位：开创性突破，还是既有纲领的凸闭包？}

\subsection{问题的提法}

一个自然的问题是：OpenAI 在 \NS{} 上的工作，究竟发明了新的数学思想与方法，还是既有数学知识与方法的闭包推论？换言之，它是否带来了数学思想层面的革命性创新？

要回答，先得把“一个结果在数学思想史中的位置”拆成可以分别判断的层次，否则“创新”一词会在不同尺度上滑动。本章采用三层判据：\textbf{纲领层}，指研究纲领、问题的重新表述与核心图景；\textbf{技术层}，指为跨越具体障碍而引入的构造步骤与估计；\textbf{方法论层}，指生产与验证知识的流程与工具。一个工作可以在某一层是革命性的，而在另一层只是收官；把三层混为一谈，是这类讨论中最常见的失误。

\subsection{纲领层：不是它发明的}

判据一：核心图景是否出自本工作。答案是否。

“先选解、后定力”的残差博弈、逐尺度放大、修正阶数随尺度上升，这三件事构成整个工作的思想内核，而它们的归属在公开材料中高度一致。Buckmaster 在声明中把纲领的基本想法记在 Córdoba 与 Martínez-Zoroa 名下，并称该攻击路线“不是由我们提出，也不是由大语言模型提出”；Tao 在 9 月 7 日（早于 OpenAI 公布结果）的博文中独立地给出同一条溯源，并明确指出从 Alpöge--Buckmaster 的强迫 Euler 到 \NS{} 的道路“对任何理解这套方法的人都是可见的”。按此，纲领层不属于 OpenAI 的原创。

需要谨慎的一点是：可预期不等于可推出。一个结果落在既有纲领的射程内，并不自动意味着它是该纲领现有定理的平凡推论——许多最终被纳入教科书的结果，同样在事前被同行预期过。因此“纲领非原创”只能否掉“新范式”这一说法，还不足以直接判定它是“凸闭包”。

\subsection{技术层：有一处实质性的新支点，但不是新范式}

判据二：跨越 \NS{} 这一具体障碍时，是否出现了既有纲领中没有、且不可由既有步骤机械合成的新构造步骤。

存在这样一处支点，也是本文认为最值得记录的地方：用零角平均振荡脉冲的 Reynolds 应力，去补足坍缩涡旋在奇点处缺失的动量输运。它的新意有两重。其一，它改变了被放大扰动的角色：在 Boussinesq 中，被放大的波本身就是奇点；在 \NS{} 中，波退居为工具，其平均自作用的散度充当作用在背景上的内部力。角色的重新指派不是尺度的微调，而是一次概念上的重定向。其二，它把“非线性项二次不消失”这一性质从障碍变成了资源：正是
\begin{equation}
\big\langle (w-\langle w\rangle)\otimes(w-\langle w\rangle)\big\rangle\neq 0
\end{equation}
这一事实，使得平均为零的扰动可以承担非零的通量。这项技术在更早的工作中已有先声（Daneri--Székelyhidi 用振荡实现指定应力、Craik--Criminale 的精确波、Córdoba--Martínez-Zoroa 的层次放大），但把“脉冲向坍缩核心补给动量”作为主线安排到 \NS{} 的强迫情形，是此处特有的组合。

与此配套的还有各向异性 Type II 坍缩：$\ell_{z}/\ell_{r}\asymp\tau^{-h}\to\infty$ 的细长柱，与以往层次构造的各向同性倾向不同。这一点的重要性后来被 CIV 的刚性定理反证式地确认——若核心不精确轴对称、或标度退化为各向同性，假想奇点将被逼回正则。因此这套形态不是任意的技术选择，而是构造得以成立的必要条件之一。

但需要同样清楚地指出：这些是技术层的支点，不是纲领层的突破；它们没有引入新的数学对象、新的不变量或新的普适结构，而是在既有纲领内部、为跨越一个具体障碍而设计的步骤。为把这一判断从断言变成可核验的结论，下面把构造的数学成分逐项拆开，核对其出处。

\subsection{成分溯源：五项可核验的出处}

\textbf{（一）迭代修正机制出自 Córdoba--Martínez-Zoroa。} 逐阶构造、把误差推向更高阶、让修正阶数随尺度上升，这些是 CMZ 纲领的规定动作。\NS{} 构造的第三至第十章是这套记账在更复杂对象上的执行，而非新机制。

\textbf{（二）精确波出自 Craik--Criminale，其不稳定性分析出自 Leibovich--Stewartson。} 第 4 章推出的振幅方程 $\dot a=-(I-2\,\xi\otimes\xi/|\xi|^{2})Da$ 与相位方程 $\dot\xi=-D^{\top}\xi$，在结构上与柱状涡的离心不稳定性分析完全一致：反射矩阵 $P_{\xi}=I-2\,\xi\otimes\xi/|\xi|^{2}$、振幅沿不稳定特征线的指数增长、波数随背景应变被拉伸，都是既有对象（见\cite{leibovich,billant}）。脉冲的生长没有引入新工具，只是把既有工具搬到环域上使用。

\textbf{（三）用振荡实现指定应力出自 Daneri--Székelyhidi，但用法不同。} 在凸积分传统中，振荡用于构造非唯一弱解、逼近一个 h-原理意义下的容许应力场，其目标落在“解的低正则自由度”上\cite{ds}。OpenAI 的用法在结构上不同：脉冲服务于一条\textbf{预先写定的光滑奇异背景}，其平均二次自作用的散度被用来抵消背景残差的奇异部分。这是“用同一材料、换一个工种”——振荡实现应力是既有的（棋类级）手段，而把应力供给绑定到一条坍缩背景上，是新的（棋步级）用法。这里随之出现一个此前没有的逆问题：给定必须实现的目标应力，构造能实现它的脉冲族；它的可解性由附录 C 的“容许应力锥”回答，而该附录在手段上仍属凸积分传统的技术（以振荡族的平均实现给定应力，可解性归结为某个锥条件），不引入新的数学对象。需要说明的是，本文未核对附录 C 的完整技术细节，这一归类以该附录的题名与相邻文献为据。

\textbf{（四）各向异性标度属于 CIV 已分析的 Type II 族，不是新标度类型。} $\ell_{r}\asymp\tau^{1/2}$、$\ell_{z}\asymp\tau^{1/2-h}$ 的奇点在轴向占优，比各向同性的 Type I（处处 $\tau^{-1/2}$）发散更快，因而落在 Type II 范畴；CIV 的刚性定理正是在这一族里工作，并证明只要再附加空间解析的外力，正则性就会回来\cite{civ}。更精确地说，CIV 在其注 1.4(b) 中明确指出，该定理与既有的一系列覆盖 Type I 情形的正则性结果——Chen--Strain--Yau--Tsai、Chen--Strain--Tsai--Yau、Koch--Nadirashvili--Seregin--Šverák、Seregin--Šverák 等——互不包含。这意味着各向异性 Type II 区间是先验定理的一块盲区：既有刚性结果既覆盖不到它、也排除不了它，而构造恰恰落在这块盲区里。因此这套标度是一个已被刻画过的族中的一点，其新颖性只在于“它是这条构造路径的必要选择”。

\textbf{（五）环面上的除子界与反演乘子是标准拟微分记账。} 引理 8.6 用到的两件工具——用除子界控制方向导数逆算子的象征、用二维 Fourier 可和性收尾——在拟微分算子与格点计数里都是常备件（第 8 章已给出完整推导）。它在这里的地位是“承重但普通”，不是新工具。

\textbf{屏障与转换：为何这一组合不可被机械合成。} 在五项溯源之外，还有一条来自方程本身的证据。CMZ 纲领的层次构造之所以能逐层升高频率，是因为在 Euler 与 IPM 中高频扰动没有黏性阻尼这个对手；而在 \NS{} 中，每个脉冲都背着 $-\nu|\xi|^{2}a$ 的衰减（第 6 章），频率越高衰减越快。这直接阻断了“把层放大原样译入 \NS{}”的机械路线，也是构造必须把扰动的角色从“奇点载体”改写为“应力供给者”的技术原因——角色转换不是风格选择，而是被黏性逼出来的。角色一旦改写，就必须回答那个此前没有的逆问题：给定必须实现的目标应力，构造能实现它的脉冲族（附录 C）。被黏性逼出的角色转换，加上随之而来的应力实现逆问题，合起来正是“不可机械合成”的具体含义。

还有一条时间线对照可以量化这道鸿沟。Albritton--Brué--Colombo\cite{abc}的非唯一性构造（2022）同样以“先选解、后定力”为骨架，其外力属于 $L_{t}^{1}L_{x}^{2}$ 且在初始时刻奇异；从这类允许奇异外力出场的结果，到备选 C 所要求的处处光滑、紧支、连同一切导数穿过奇点的外力，中间隔着的正是脉冲机制所要弥合的那段距离，也是这条路线此后数年主攻的方向。

把五项溯源与上述屏障证据并起来，得到一个比断言更强的判断：构造所用的每一块数学材料，都能在 2026 年之前已存在的文献里找到出处，没有任何一块是为此新造的数学对象、新不变量或新结构；而把这些既有材料拼装为有效证明的那一步——角色转换与应力实现——是可指认的具体新支点。这正是“它发明了一步新棋，而不是一种新棋类”的技术含义。

\subsection{方法论层：真正的新在流程，不在数学}

判据三：生产与验证知识的流程是否发生了阶跃。在这里，答案倾向于肯定，但必须明确这一肯定不落在数学思想层。

把一份由模型生成、据其共同作者自述“差到难以卒读”的写稿，在约一周内变成可被 Lean 内核编译的证书；把约一万个智能体组织成并行探索不同问题变体的搜索系统；用一次大模型调用完成对海量代码的形式化与验证——这些是若干年前不可想象的流程能力，它们改变的是求解与验证的成本结构。但按本层的判据，这属于方法论与工程，而不属于数学思想本身；把它记入“数学革命”的名下，会混淆两类不同的创新。

\subsection{综合判断}

把三层并起来，得到一个尽量不含混的定位。

在纲领层，它是既有纲领（Córdoba--Martínez-Zoroa 的残差博弈与逐尺度放大）在 \NS{} 强迫情形上的收官，思想内核非其原创。在技术层，它含有一处实质性的新支点（脉冲经 Reynolds 应力补缺力）与一套被刚性定理确认必要的各向异性坍缩形态，足以使它不落为机械合成，但尚不构成新范式。在方法论层，它是真正的新东西，也是本次事件中最可能被后人反复提及的部分，但它属于流程而非数学。

因此，若把“革命性的数学思想创新”定义为引入新对象、新结构或新不变量，并据此改写领域的问题框架，那么这项工作不是。若把它定义为“在既有纲领内，以一个此前没有的新支点跨过不可逾越的具体障碍”，那么它是，并且是其中规格相当高的一次。

\begin{table}[H]
\centering
\small
\caption{“思想定位”可操作判据表}
\begin{tabular}{@{}p{5.2cm}p{3.2cm}p{5.6cm}@{}}
\toprule
\hdrow \textbf{判据} & \textbf{本工作的答案} & \textbf{依据} \\
\midrule
是否引入新的数学对象或不变量 & 否 & 沿 Fefferman 既有提法与既有方程 \\
是否改变领域的问题框架 & 否 & 求解的是既有的备选 C、D \\
是否提出纲领之外的新原理 & 否 & Córdoba--Martínez-Zoroa 的残差博弈与逐尺度放大 \\
是否提供不可机械合成的关键步骤 & 是 & 脉冲经平均二次自作用补足坍缩缺力 \\
该步骤是否被后续工作确认为必要 & 是 & CIV 刚性定理排除各向同性与解析外力 \\
是否带来生产与验证流程的阶跃 & 是 & 多智能体搜索与规模化自动形式化（属方法论） \\
\bottomrule
\end{tabular}
\end{table}

最后需要一层诚实：上述判定以“我们确实理解了这份证明的动作”为前提，而这一点目前并不完全成立。形式化证书是自评的、尚未被独立核验；语义忠实性存在具体争议（第 8、9 章）；数学思想尚未被共同体严格重述。在证明被真正消化之前，“它在思想史中属于开创还是收官”这个问题只能给出暂定答案。这不是回避，而是对判据适用条件的如实交代。

\section{结语：修正后的判断}

本文开始时提出了一个问题：如何定位 OpenAI 在 \NS{} 上的这项工作。经过前十一章的梳理，可以给出一个尽量收敛、同时保留边界的判断。

\textbf{第一，它首先是一条技术路线的收敛证据。} 从 Córdoba--Martínez-Zoroa 的残差博弈，到 Alpöge--Buckmaster 的光滑外力与 Euler，再到把 \NS{} 的坍缩核心与由离心不稳定性生长的脉冲对接起来，这条路线有一以贯之的思想内核：先写下想要的奇异流，再让外力等于残差，然后集中全部技巧去保证残差光滑。OpenAI 的贡献在于把这套思想推到了 \NS{} 的强迫情形，其独有的一步是用零角平均脉冲的 Reynolds 应力补足坍缩缺失的动量输运。

\textbf{第二，它同时是一次对形式化验证认识的施压。} “Lean 通过”与“证明正确”之间，隔着“语义忠实翻译”这道在 SCI 谱系中高达无穷的坎；引理 8.6 处 $m+4$ 与 $m+5$ 的差异，是这道坎的一个具体样本。这条教训不针对 OpenAI，而适用于一切以自动形式化作为验证手段的工作。

\textbf{第三，它不是千禧年问题的终局裁判。} 备选 A 与 B 仍然开放；“有外力则可爆破”与“无外力是否永远光滑”本是两问。把它读成“AI 证明了水的方程会爆炸”，是对官方提法的误读。

\textbf{第四，边界必须写清楚。} 本文未独立复核该证明；关于正确性的判断均转述自公开文献并注明立场；CIV 的刚性定理是约束而非否证；争议事实均逐项标明来源、按归因分别叙述。这些限制不是谦辞，而是这篇综述得以成立的诚实前提。

把这四条并起来，得到的定位是：这是一次真实的、可检验的数学推进，也是一次对“什么才算验证”的追问。它推进的是一种构造方法，追问的是一种信念——我们是否愿意接受一个正确但尚未被理解的证明，作为数学知识被共同持有。

\begin{keybox}
\textbf{一句话收束}\quad 这项工作把一条酝酿多年的构造路线推到了 \NS{} 的强迫情形，也把“形式化通过是否等于理解发生”这个问题，第一次以如此具体的形态摊在数学共同体面前。前者是数学，后者是元数学，都还没有结束。
\end{keybox}

\appendix

\section{记号速查表}

\begin{table}[H]
\centering
\small
\caption{全文主要记号}
\begin{tabular}{@{}p{2.4cm}p{11.6cm}@{}}
\toprule
\hdrow \textbf{记号} & \textbf{含义} \\
\midrule
$u,p,f,\nu$ & 速度、压力、外力、黏性系数 \\
$\tau$ & 奇点前的剩余时间 $1-t$ \\
$\ell_{r},\ell_{z}$ & 核心的径向、轴向特征长度 \\
$h$ & 各向异性指数，$0<h<1/100$ \\
$\mathcal{P}u$ & 速度关于 $z$ 轴的角向平均（轴对称部分） \\
$w=u-\mathcal{P}u$ & 速度的非轴对称部分 \\
$\langle\cdot\rangle$ & 角向平均算子 \\
A/B/C/D & 千禧年问题官方提法的四条备选 \\
$\mathbb{T}^{3},\mathbb{T}^{2}$ & 三维、二维环面 \\
SCI & 可解性复杂度层级（Solvability Complexity Index） \\
\bottomrule
\end{tabular}
\end{table}

\section{一致性元数据}

\begin{table}[H]
\centering
\footnotesize
\caption{关键论断与可核验来源的对应}
\begin{tabular}{@{}p{5.2cm}p{8.8cm}@{}}
\toprule
\hdrow \textbf{本文论断} & \textbf{可核验来源} \\
\midrule
官方四备选 A--D 的表述 & Fefferman 撰写的 Clay 官方提法\cite{clay} \\
定理 1.1（全空间）与推论 10.6（周期） & OpenAI 手稿\cite{openai-paper} 第 1 节、第 10 节 \\
约一万并发智能体、88 小时、约 1300 亿输出 token & OpenAI 官方说明\cite{openai-blog} \\
Lean 仓库评审状态为“自评”、无 \texttt{sorry} & OpenAI Lean 仓库\cite{openai-lean}；Stanford Tech Review 审计\cite{stanford-audit} \\
Córdoba--Martínez-Zoroa 的 IPM 爆破 & \arxiv{2410.22920}\cite{cmz-ipm} \\
Alpöge--Buckmaster 的归属声明 & Buckmaster 公开声明\cite{buckmaster-statement} \\
各向异性标度 $\ell_{r}\asymp\tau^{1/2}$、$\ell_{z}\asymp\tau^{1/2-h}$ & OpenAI 手稿\cite{openai-paper} 第 2 节；CIV\cite{civ} 第 1 节 \\
引理 8.6 的 $m+4$ 与 Lean 的 $m+5$ & Bastounis--Circelli--Hansen\cite{bch} 例 3.1；并见本文据 OpenAI Lean 仓库\cite{openai-lean} commit \texttt{f9e8bc5} 的一手核验（第 8.3 节） \\
刚性定理（Type II ＋轴对称核心 ＋解析外力 $\Rightarrow$ 正则） & Constantin--Ignatova--Vicol\cite{civ} 定理 1.1 及其推论 \\
Clay 的“刻意不匆忙”表态 & Clay 新闻页\cite{clay-news} \\
正缺陷问题作为无外力靶标 & Petrillo--Glimm\arxiv{2609.23868}\cite{petrillo-glimm} \\
五个数学成分的出处（迭代修正、精确波、振荡实现应力、Type II 标度、环面乘子）与“屏障与转换”论证 & CMZ\cite{cmz-ipm}；Craik--Criminale 与 Leibovich--Stewartson\cite{leibovich,billant}；Daneri--Sz\'ekelyhidi\cite{ds}；CIV\cite{civ}；Albritton--Bru\'e--Colombo\cite{abc} \\
\bottomrule
\end{tabular}
\end{table}

\clearpage

\begin{thebibliography}{99}
\small
\setlength{\itemsep}{2pt}

\bibitem{clay} Fefferman C L. Existence and smoothness of the Navier--Stokes equation[C/OL]//The Millennium Prize Problems. Clay Mathematics Institute, 2000. \url{https://www.claymath.org/wp-content/uploads/2022/06/navierstokes.pdf}.

\bibitem{openai-blog} OpenAI. On the Navier--Stokes Millennium Prize Problem[EB/OL]. 2026-09-08. \url{https://openai.com/index/navier-stokes-solution/}.

\bibitem{openai-paper} OpenAI. Finite time blowup for Navier--Stokes[Z/OL]. 手稿, 2026. \url{https://cdn.openai.com/pdf/32d9f210-8b73-45e0-91bc-82a30aef8a9a/navier-stokes.pdf}.

\bibitem{openai-lean} OpenAI. NavierStokesAndEuler: Lean 4 formalizations[CP/OL]. GitHub. \url{https://github.com/openai/NavierStokesAndEuler}.

\bibitem{clay-news} Clay Mathematics Institute. Navier--Stokes announcement[EB/OL]. 2026-09-11. \url{https://www.claymath.org/news/navier-stokes-announcement/}.

\bibitem{leray} Leray J. Sur le mouvement d'un liquide visqueux emplissant l'espace[J]. Acta Mathematica, 1934, 63: 193--248.

\bibitem{ckn} Caffarelli L, Kohn R, Nirenberg L. Partial regularity of suitable weak solutions of the Navier--Stokes equations[J]. Communications on Pure and Applied Mathematics, 1982, 35(6): 771--831.

\bibitem{ess} Escauriaza L, Seregin G, \v{S}ver\'ak V. $L_{3,\infty}$-solutions of Navier--Stokes equations and backward uniqueness[J]. Russian Mathematical Surveys, 2003, 58(2): 211--250.

\bibitem{tao-averaged} Tao T. Finite time blowup for an averaged three-dimensional Navier--Stokes equation[J]. Journal of the American Mathematical Society, 2016, 29(3): 601--674.

\bibitem{bv-nonunique} Buckmaster T, Vicol V. Nonuniqueness of weak solutions to the Navier--Stokes equation[J]. Annals of Mathematics, 2019, 189(1): 101--144.

\bibitem{abc} Albritton D, Bru\'e E, Colombo M. Non-uniqueness of Leray solutions of the forced Navier--Stokes equations[J]. Annals of Mathematics, 2022, 196(1): 415--455.

\bibitem{ds} Daneri S, Sz\'ekelyhidi Jr L. Non-uniqueness and h-principle for H\"older-continuous weak solutions of the Euler equations[J]. Archive for Rational Mechanics and Analysis, 2017, 224(2): 471--514. DOI: 10.1007/s00205-017-1089-0.

\bibitem{cmz-ipm} C\'ordoba D, Mart\'inez-Zoroa L. Finite time singularities of smooth solutions for the 2D incompressible porous media (IPM) equation with a smooth source[Z/OL]. \arxiv{2410.22920}, 2024.

\bibitem{cmz-euler} C\'ordoba D, Mart\'inez-Zoroa L, Zheng F. 光滑外力下流体方程奇点形成的相关工作[Z]. 见手稿\cite{openai-paper} 参考文献 [6][8]。

\bibitem{buckmaster-statement} Buckmaster T. Statement[EB/OL]. 2026-09-08. \url{https://cims.nyu.edu/~tristanb/statement.pdf}.

\bibitem{tao-blog-sep7} Tao T. Finite time blowup with smooth forcing term for the incompressible porous medium, Boussinesq, and incompressible Euler equations[EB/OL]. What's new, 2026-09-07. \url{https://terrytao.wordpress.com/2026/09/07/}.

\bibitem{schneider} Schneider T. Headlines and inside stories: understanding and trust in AI for mathematics, science, and engineering[EB/OL]. 客座于 What's new, 2026-09-23. \url{https://terrytao.wordpress.com/2026/09/23/}.

\bibitem{gowers} Gowers T. Why I didn't sign the Fields medallists' letter[EB/OL]. 2026-09-17. \url{https://terrytao.wordpress.com/2026/09/17/}.

\bibitem{mathandai} 25 Fields medallists. Letter on the misalignment between AI companies and the mathematical community[EB/OL]. \url{https://mathandai.org/}.

\bibitem{civ} Constantin P, Ignatova M, Vicol V. Regularity of asymptotically axisymmetric solutions to the 3D Navier--Stokes equations with analytic forcing[Z/OL]. \arxiv{2609.20803}, 2026.

\bibitem{bch} Bastounis A, Circelli F, Hansen A C. Navier--Stokes lost in translation: why Lean verification of AI autoformalisation does not guarantee correct natural language proofs[Z/OL]. \arxiv{2610.08144}, 2026.

\bibitem{petrillo-glimm} Petrillo J, Glimm J. The positive defect problem: target and admissibility criteria for a programmatic search for unforced Navier--Stokes blowup[Z/OL]. \arxiv{2609.23868}, 2026.

\bibitem{shavit} Shavit M. Two forced finite-time singularities, and the mechanism behind each: a friendly reading of Alp\"oge--Buckmaster (Boussinesq) and OpenAI (Navier--Stokes)[EB/OL]. 2026-09. \url{https://michalshavitnyu.github.io/files/blowup-mechanisms.pdf}.

\bibitem{stanford-audit} Raman P. OpenAI vs Buckmaster: the Navier--Stokes Lean proofs, audited[EB/OL]. Stanford Tech Review, 2026-09-08. \url{https://stanfordtechreview.com/articles/openai-buckmaster-navier-stokes-lean-proofs}.

\bibitem{leibovich} Leibovich S, Stewartson K. A sufficient condition for the instability of columnar vortices[J]. Journal of Fluid Mechanics, 1983, 126: 335--356.

\bibitem{billant} Billant P, Gallaire F. Generalized Rayleigh criterion for non-axisymmetric centrifugal instabilities[J]. Journal of Fluid Mechanics, 2005, 542: 365--379.

\end{thebibliography}